\documentclass[11pt]{article}
\usepackage[a4paper,margin=25mm]{geometry}
\usepackage{amsmath,amssymb}
\usepackage[T1]{fontenc}
\setlength{\emergencystretch}{2em}
\title{Proof of conjecture 00000006340}
\author{ziangni-sys}
\date{}
\begin{document}
\maketitle
\noindent\textbf{Source and scope.}
Both source versions ask for an explicit orthogonal pair of matrices with the same rigidity function (the same rank-change costs) but different eigenvalues. We give such a pair using the standard entry-change matrix rigidity. Spectral rigidity is not assigned a separate formula in the source; our separation proves its stated concrete realization by distinct spectra. The pair consists of orthogonal matrices and is also orthogonal in the Frobenius inner product.

\medskip
\noindent\textbf{The matrices.}
Set
\[
 A=I_2=\begin{pmatrix}1&0\\0&1\end{pmatrix},
 \qquad J=\begin{pmatrix}0&-1\\1&0\end{pmatrix},
 \qquad K=J^T=\begin{pmatrix}0&1\\-1&0\end{pmatrix}.
\]
Direct multiplication gives
\[
 J^TJ=JJ^T=I_2,\quad KJ=JK=I_2,\quad
 A^TA=AA^T=I_2.
\]
Furthermore $\langle A,J\rangle_F=\sum_{i,j}A_{ij}J_{ij}=0$.
Both meanings of an orthogonal pair, individually orthogonal matrices and Frobenius orthogonality, hold.

\medskip
\noindent\textbf{The genuine rigidity function.}
For real $2\times2$ matrices $M,N$, let
\[
 c(M,N)=\#\{(i,j):M_{ij}\ne N_{ij}\},\qquad
 R_M(r)=\min\{c(M,N):\operatorname{rank}N\le r\},
 \quad r\in\mathbb N.
\]
The feasible set is nonempty, since $N=0$ is allowed. The costs are natural numbers, so the minimum exists and is attained. This is the usual rigidity function: the fewest entries one must change to reduce the rank to at most $r$. Equivalently one may write $N=M+E$ and count nonzero entries of $E$.

\medskip
\noindent\textbf{An exact bijection of all perturbations.}
Left multiplication by $J$ is a signed row permutation:
\[
 JN=\begin{pmatrix}-N_{21}&-N_{22}\\N_{11}&N_{12}\end{pmatrix}.
\]
It is invertible, so $\operatorname{rank}(JN)=\operatorname{rank}N$.
For any $M,N$, two entries are unequal exactly when their simultaneous negatives are unequal; a row permutation merely permutes these tests. Consequently
\[
 c(JM,JN)=c(M,N).
\]
In particular, $JA=J$ and
\[
 c(J,JN)=c(A,N).
\]
The map $N\mapsto JN$ therefore sends every rank-$\le r$ perturbation of $A$ to a rank-$\le r$ perturbation of $J$ with the identical cost. Its inverse is $N\mapsto KN$, which also preserves rank and changed-entry counts. Hence the complete sets of attainable costs agree for every $r$, and
\[
 \boxed{R_A(r)=R_J(r)\quad\text{for every }r\in\mathbb N.}
\]
This proves equality of the full rigidity functions, not just one selected rank or a common upper bound.

\newpage
\noindent\textbf{The actual eigenvalue separation.}
Complex eigenvalues are characterized by actual nonzero eigenvectors:
$\lambda\in\sigma(M)$ if and only if $Mv=\lambda v$ for some
$v\in\mathbb C^2\setminus\{0\}$. For $A=I_2$, this gives
$\sigma(A)=\{1\}$ immediately.

For $J$, the eigenvector equations are
\[
 -v_2=\lambda v_1,\qquad v_1=\lambda v_2.
\]
A nonzero eigenvector must have $v_1\ne0$: otherwise the first equation also forces $v_2=0$. Combining the equations gives
$(\lambda^2+1)v_1=0$, hence $\lambda^2+1=0$ and
$\lambda=i$ or $\lambda=-i$. Conversely,
\[
 J\binom{1}{-i}=i\binom{1}{-i},
 \qquad
 J\binom{1}{i}=-i\binom{1}{i}.
\]
Both vectors are nonzero. Thus
\[
 \boxed{\sigma(A)=\{1\},\qquad \sigma(J)=\{i,-i\}.}
\]
In particular $1$ belongs to the first spectrum and not the second, proving that the spectra differ.

\medskip
\noindent\textbf{Formal correspondence.}
The Lean development uses actual real matrices and Mathlib's matrix rank, defined as the finite dimension of the range of the associated linear map. The changed-entry count is the finite sum of the zero--one indicator of inequality at each entry. Actual natural-number cost sets define the rigidity as their \texttt{sInf}; nonemptiness and the well-ordering of the natural numbers prove this infimum is an attained minimum and is at most every feasible cost.

The determinant identities $\det J=\det K=1$ prove rank preservation for every matrix through Mathlib's invertible-multiplication rank theorem. Explicit matrix multiplication proves entry-count preservation. Both directions of the perturbation transformation prove equality of the entire attainable-cost sets, and therefore equality of the entire rigidity functions.

The source matrices are real, and their complexification is the actual entrywise inclusion into $\mathbb C$. The eigenvalue definition quantifies over actual nonzero complex vectors and the matrix-vector equation. The formal proof characterizes all eigenvalues in both directions, including explicit witnesses for both quarter-turn eigenvalues; it does not merely supply possible eigenvalues or characteristic-polynomial guesses. Matrix orthogonality and Frobenius orthogonality are separately verified.

\medskip
\noindent\textbf{Reproduction and validation.}
Run \texttt{lake build} in the submitted \texttt{lean} directory with Lean 4.19.0. Mathlib is pinned publicly to
\[
 \texttt{c44e0c8ee63ca166450922a373c7409c5d26b00b}.
\]
Printed axiom audits cover orthogonality, rank and cost preservation, attainable-cost equality, the genuine minimum, the full rigidity equality, both eigenvalue characterizations, the actual spectra and the final separation. The submission includes the complete LaTeX source and its compiled PDF.

\medskip
\noindent\textbf{Conclusion.}
The explicit identity and quarter-turn matrices form the requested orthogonal pair: all rank-change costs, and thus the full standard rigidity functions, agree, while the eigenvalues differ. This establishes the concrete separation requested in both source versions.
\end{document}
